\section{Candidate filtering and independent schedule replay}
\label{app:implementation}

\subsection{Time-support filtering of cargo variables}
\label{app:filter}

The sparse variant removes cargo variables, while keeping the vehicle and
wave-selector mother domains. It therefore preserves automated repositioning,
independent empty transfers, and zero-cargo vehicle movements. The filter
does not infer future city availability from historical idle stock alone:
a vehicle that is absent now may arrive before a future task.

For a revealed order $\ell\in\mathcal L^d\cap\mathcal L(t)$, form provisional new
collection candidates $P_\ell$ obeying \eqref{alg:origin-cut},
$i\geq s_\ell$, and positive local capacity at the upper fleet size $N$.
Form provisional new distribution candidates $R_\ell$ obeying
\eqref{alg:downstream-cut}, $j\leq e_\ell$, and positive local capacity.
Let $J^P_\ell$ be the end times of committed positive collection cargo and
$J^A_\ell$ the arrival times of committed positive automated cargo. These
history sets seed downstream possibilities even when there is no admissible
new collection tour. Define
\begin{equation}
 \underline t_\ell^P=\min\bigl(\{j:(i,j)\in P_\ell\}\cup J^P_\ell\bigr),
 \qquad
 \overline t_\ell^R=\max\{i:(i,j)\in R_\ell\},
 \label{alg:time-bounds}
\end{equation}
with the empty-set values $+\infty$ and $-\infty$, respectively. If
$A>0$, retain provisional automated cargo arcs
\begin{equation}
 Q_\ell=\{(i,j)\in\mathcal A_A:
 t\leq i<c_t,\ j\leq c_t,\quad
 \underline t_\ell^P\leq i,\ j\leq\overline t_\ell^R\};
 \label{alg:auto-support}
\end{equation}
if $A=0$, take $Q_\ell=\varnothing$. A new collection cargo variable is
retained only if its end is no later than some start in $Q_\ell$.
A new distribution cargo variable is retained only if its start is at least
\begin{equation}
 \min\bigl(\{j:(i,j)\in Q_\ell\}\cup J^A_\ell\bigr).
 \label{alg:delivery-support}
\end{equation}
Direct cargo candidates must obey $i\geq\max\{t,s_\ell\}$,
$i<b_t$, $j\leq\min\{c_t,e_\ell\}$, and
$\widetilde K^d_{ij}(N)>0$. If $N=0$, all new human-driven cargo
candidates are empty. All previously committed nonzero cargo keys are
retained regardless of these tests so that history can be fixed exactly.

These tests use time support and optimistic capacity; they do not solve
the joint fleet-allocation problem. Retaining a candidate is therefore no
claim that the candidate is feasible together with other decisions. In
the MILP, omitted cargo variables are zero, and every remaining cargo
variable is still subject to the original capacity, vehicle, precedence,
and demand constraints. The procedure records the number of cargo
variables before and after filtering for each rebuilt window. We do not
assert a general preservation theorem or a computational speed advantage
from this filter; neither follows from the construction alone.

\subsection{Independent replay and certificate scope}
\label{app:replay}

A returned schedule is checked independently of the optimizer's capacity
tables. For example, a positive local wave load $q$ and count $v$ must
satisfy
\begin{equation}
 q\leq vM,\qquad
 f_k(q)\leq v[(j-i)\Delta t-2\rho_k]^+,
 \label{alg:local-replay}
\end{equation}
and positive direct cargo must satisfy $q\leq vM$ and
\eqref{mod:direct-duration} within its assigned duration. Replay also
checks domains, selector links, shared automated capacity, every city
inventory, order windows, cumulative precedence, and demand balance.
Costs are recomputed from initiated whole tasks once, rather than by
summing overlapping window objective values.

For a completed offline solve, a reported MILP gap refers only to the
specified finite-horizon model and its input information. Window solver
bounds and gaps refer to their individual restricted subproblems; the
last window's bound is not a global lower bound for rolling execution.
For partial execution, replay can certify past commitments and accounting,
including hub or in-transit goods, but defers final stage balance and future
completion reachability. A partial schedule consequently has no certified
complete-horizon objective. These distinctions are also needed when
comparing this finite-fleet dispatch procedure with the reduced
steady-state and QVI models.

